\documentclass[12pt,letterpaper]{article}

\usepackage[margin=3cm]{geometry}
%\usepackage{chemfig}

\begin{document}

% This is a comment and will not be printed

One of 
the main appeals 
of \LaTeX\ is that it is \textit{very} good at
\textbf{typesetting} equations:

$$
a_n = \left(\frac{1}{\pi}\right)\cos(4x)
$$

\begin{equation}
\nabla \times \mathbf {E} =-{\frac {1}{c}}{\frac {\partial \mathbf {B} }{\partial t}}
\label{my_equation}
\end{equation}



\begin{equation}
\theta =\theta _{0}+\omega t-{\frac{1}{2}}\alpha t^{2}
\label{my_other_equation}
\end{equation}

The advantage of the \texttt{equation} environment is that is allows 
you to label and reference equations like you see above in Equations 
\ref{my_equation} and \ref{my_other_equation}.

\clearpage

This is another page

\end{document}

Nothing below the end command will be processed. I find this a useful place to put drafts/notes and stuff

\begin{equation}
|\psi \rangle =a_{\psi }|\uparrow _{z}\rangle +b_{\psi }|\downarrow _{z}\rangle
\label{my_other_other_equation}
\end{equation}

\vspace{2cm}
\chemfig{H-C(-[2]H)(-[6]H)-C(=[1]O)-[7]H}
